\documentclass[9pt,a4paper]{extarticle}

% ---------- Compact ICPC-style Python Team Reference ----------
\usepackage[a4paper,top=0.50in,bottom=0.45in,left=0.45in,right=0.45in]{geometry}
\usepackage[T1]{fontenc}
\usepackage{lmodern}
\usepackage{amsmath,amssymb}
\usepackage{multicol}
\usepackage{listings}
\usepackage{xcolor}
\usepackage{fancyhdr}
\usepackage{enumitem}
\usepackage{titlesec}
\usepackage{microtype}
\usepackage{transparent}

% Preserve leading spaces better when copying code from a pdfLaTeX PDF.
\makeatletter
\def\@xobeysp{\leavevmode\nobreak\texttransparent{0}{\char32}}
\makeatother

% ---------- Team information ----------
\newcommand{\UniversityName}{United International University}
\newcommand{\TeamName}{Team: \rule{3.0cm}{0.4pt}}

% ---------- Header / footer ----------
\pagestyle{fancy}
\fancyhf{}
\fancyhead[L]{\footnotesize \UniversityName\ --- \TeamName}
\fancyhead[R]{\footnotesize Page \thepage}
\renewcommand{\headrulewidth}{0.3pt}
\setlength{\headheight}{12pt}

% ---------- Spacing ----------
\setlength{\parindent}{0pt}
\setlength{\parskip}{1pt}
\setlength{\columnsep}{0.18in}
\setlength{\columnseprule}{0.15pt}
\setlist{nosep,leftmargin=*}

\titleformat{\section}
  {\normalsize\bfseries}
  {}{0pt}{}
  [\vspace{-2pt}\hrule]
\titleformat{\subsection}
  {\small\bfseries}
  {}{0pt}{}
\titlespacing*{\section}{0pt}{4pt}{3pt}
\titlespacing*{\subsection}{0pt}{3pt}{1pt}

% ---------- Code style: monochrome to save ink ----------
\lstset{
  language=Python,
  basicstyle=\ttfamily\footnotesize,
  keywordstyle=\bfseries,
  commentstyle=\itshape,
  stringstyle=\ttfamily,
  showstringspaces=false,
  keepspaces=true,
  columns=flexible,
  breaklines=false,
  tabsize=4,
  frame=none,
  aboveskip=1pt,
  belowskip=2pt,
  xleftmargin=0pt,
  xrightmargin=0pt
}

\newcommand{\code}[1]{\texttt{#1}}

\begin{document}
\small

\begin{center}
{\large\bfseries Python Competitive Programming Reference}\\[-2pt]
{\footnotesize Compact team reference --- algorithms, mathematics, graphs, and data structures}
\end{center}
\vspace{-4pt}

\begin{multicols}{2}

% ============================================================
\section{Python Syntax, Formatting \& Fast I/O}

\begin{lstlisting}
# Binary string -> integer
n = int(s, 2)

# Integer -> zero-padded binary string
b = format(x, f"0{width}b")

# Float formatting
print("%.2f" % value)
print(f"{answer:.2f}")

# Read entire input as integers
import sys
data = list(map(int, sys.stdin.buffer.read().split()))

# Read entire input as strings
words = sys.stdin.buffer.read().split()
# bytes -> str only when needed: token.decode()
\end{lstlisting}

\subsection{Sorting a \texttt{defaultdict}}
\begin{lstlisting}
from collections import defaultdict

# By key
D = defaultdict(d.default_factory,
                sorted(d.items()))

# By value
D = defaultdict(d.default_factory,
                sorted(d.items(),
                       key=lambda kv: kv[1]))
\end{lstlisting}


% ============================================================
\section{Modular Arithmetic}

For modulus $M$:
\[
(a+b)\bmod M=((a\bmod M)+(b\bmod M))\bmod M
\]
\[
(a-b)\bmod M=((a\bmod M)-(b\bmod M)+M)\bmod M
\]
\[
(ab)\bmod M=((a\bmod M)(b\bmod M))\bmod M
\]

In Python, \code{a \% M} is already non-negative for positive $M$.

% ============================================================
\section{Sequences \& Series}

\begin{align*}
\sum_{i=1}^{n} i &= \frac{n(n+1)}2 \\
\sum_{i=1}^{n} i^2 &= \frac{n(n+1)(2n+1)}6 \\
\sum_{i=1}^{n} i^3 &= \left(\frac{n(n+1)}2\right)^2 \\
\sum_{k=0}^{n-1}(a+kd) &= \frac{n(2a+(n-1)d)}2 \\
\sum_{i=1}^{n}(2i-1) &= n^2 \\
\sum_{i=1}^{n}2i &= n(n+1)
\end{align*}

Arithmetic progression squares:
\begin{align*}
\sum_{k=0}^{n-1}(a+kd)^2
={}&na^2+a d\,n(n-1)\\
&+d^2\frac{n(n-1)(2n-1)}6
\end{align*}

Arithmetic progression cubes:
\begin{align*}
\sum_{k=0}^{n-1}(a+kd)^3
={}&na^3
+3a^2d\frac{n(n-1)}2\\
&+3ad^2\frac{n(n-1)(2n-1)}6\\
&+d^3\left(\frac{n(n-1)}2\right)^2
\end{align*}

Geometric progression ($r\ne1$):
\[
\sum_{k=0}^{n-1} ar^k=a\frac{r^n-1}{r-1}
\]
For $r=1$, the sum is $an$.

\subsection{Alternating Sum}
\begin{lstlisting}
# 1 - 2 + 3 - 4 + ... +/- n
alt_sum = -n // 2 if n % 2 == 0 else (n + 1) // 2

# 1^k - 2^k + 3^k - ... +/- n^k
alt_pow_sum = sum((1 if i & 1 else -1) * i**k
                  for i in range(1, n + 1))
\end{lstlisting}

% ============================================================
\section{Binary Search on Answer}

\subsection{First True / Minimum Valid: FFFFFTTTTT}
\begin{lstlisting}
def check(x):
    return ...

l, r = ..., ...
ans = -1
while l <= r:
    m = (l + r) // 2
    if check(m):
        ans = m
        r = m - 1
    else:
        l = m + 1
print(ans)
\end{lstlisting}

\subsection{Last True / Maximum Valid: TTTTTFFFFF}
\begin{lstlisting}
def check(x):
    return ...

l, r = ..., ...
ans = -1
while l <= r:
    m = (l + r) // 2
    if check(m):
        ans = m
        l = m + 1
    else:
        r = m - 1
print(ans)
\end{lstlisting}

\textbf{Requirement:} \code{check(x)} must be monotonic over the search interval.

% ============================================================
\section{Number Theory}

\subsection{GCD / LCM}
\begin{lstlisting}
from math import gcd
lcm = a // gcd(a, b) * b

# Euclid; useful when b itself is supplied via b % a
# without constructing a huge decimal integer.
def gcd_euclid(a, b):
    while a:
        a, b = b % a, a
    return b
\end{lstlisting}

\subsection{Prime Factorization by Trial Division}
\begin{lstlisting}
from collections import defaultdict
from math import isqrt

def factors(n):
    f = defaultdict(int)
    while n % 2 == 0:
        f[2] += 1
        n //= 2
    p = 3
    while p * p <= n:
        while n % p == 0:
            f[p] += 1
            n //= p
        p += 2
    if n > 1:
        f[n] += 1
    return f
\end{lstlisting}

If $n=\prod p_i^{e_i}$:
\[
\tau(n)=\prod(e_i+1),\qquad
\sigma(n)=\prod\frac{p_i^{e_i+1}-1}{p_i-1}.
\]

\begin{lstlisting}
def num_divisors(fact):
    ans = 1
    for e in fact.values():
        ans *= e + 1
    return ans
\end{lstlisting}

\subsection{Divisor Count / Sum for Every $1\ldots N$}
\begin{lstlisting}
N = 10**5
divcnt = [0] * (N + 1)
for i in range(1, N + 1):
    for j in range(i, N + 1, i):
        divcnt[j] += 1

# Sum of divisors instead:
# divsum[j] += i
\end{lstlisting}

\subsection{Ordinary Sieve of Eratosthenes}
\begin{lstlisting}
def sieve(N):
    is_prime = bytearray(b'\x01') * (N + 1)
    if N >= 0: is_prime[0] = 0
    if N >= 1: is_prime[1] = 0
    for p in range(2, int(N**0.5) + 1):
        if is_prime[p]:
            start = p * p
            is_prime[start:N+1:p] = b'\x00' * (
                (N - start) // p + 1)
    return is_prime
\end{lstlisting}

\subsection{Odd-Only Bytearray Sieve}
Index $i$ represents the odd number $2i+1$.
\begin{lstlisting}
from math import isqrt

def sieve_bytearray(N):
    size = (N + 1) // 2
    s = bytearray(b'\x01') * size
    if size:
        s[0] = 0                  # 1 is not prime

    lim = (isqrt(N) - 1) // 2
    for i in range(1, lim + 1):
        if s[i]:
            p = 2 * i + 1
            start = (p * p) // 2
            s[start::p] = b'\x00' * (
                (size - start - 1) // p + 1)
    return s

N = 10**6 + 7
a = sieve_bytearray(N)

def is_prime_fast(x):
    if x == 2:
        return True
    if x < 2 or x % 2 == 0 or x > N:
        return False
    return bool(a[x // 2])
\end{lstlisting}

\subsection{Compact Prime List from Odd Sieve}
\begin{lstlisting}
from array import array

primes = array('I', [2])
for i in range(1, len(a)):
    if a[i]:
        primes.append(2 * i + 1)
\end{lstlisting}

\subsection{Smallest Prime Factor (SPF) --- Many Queries}
Preprocess once; each factorization takes $O(\log n)$ factors with multiplicity.
\begin{lstlisting}
N = 10**6 + 9
spf = list(range(N + 1))

for i in range(2, int(N**0.5) + 1):
    if spf[i] == i:
        for j in range(i * i, N + 1, i):
            if spf[j] == j:
                spf[j] = i

def factorize(n):
    res = []
    while n > 1:
        p = spf[n]
        res.append(p)
        n //= p
    return res
\end{lstlisting}

\subsection{All Divisors of One Number}
\begin{lstlisting}
def all_divisors(n):
    d = []
    i = 1
    while i * i <= n:
        if n % i == 0:
            d.append(i)
            if i != n // i:
                d.append(n // i)
        i += 1
    return d              # sort(d) if order matters
\end{lstlisting}

\subsection{Goldbach Pair for an Even Integer}
\begin{lstlisting}
def goldbach(n, is_prime):
    if n < 4 or n % 2:
        return None
    for p in range(2, n // 2 + 1):
        if is_prime[p] and is_prime[n - p]:
            return p, n - p
    return None
\end{lstlisting}

\subsection{Binary Exponentiation}
\begin{lstlisting}
def bin_exp(a, b, mod=None):
    res = 1
    while b:
        if b & 1:
            res = res * a if mod is None else res * a % mod
        a = a * a if mod is None else a * a % mod
        b >>= 1
    return res

# Python built-in modular exponentiation:
# pow(a, b, MOD)
\end{lstlisting}

\subsection{Deterministic Miller--Rabin for 64-bit Integers}
\begin{lstlisting}
def is_prime_mr(n):
    if n < 2:
        return False
    small = (2, 3, 5, 7, 11, 13, 17, 19, 23, 29, 31, 37)
    for p in small:
        if n % p == 0:
            return n == p

    d = n - 1
    s = 0
    while d % 2 == 0:
        d //= 2
        s += 1

    for a in (2, 325, 9375, 28178, 450775,
              9780504, 1795265022):
        if a % n == 0:
            continue
        x = pow(a, d, n)
        if x == 1 or x == n - 1:
            continue
        for _ in range(s - 1):
            x = x * x % n
            if x == n - 1:
                break
        else:
            return False
    return True
\end{lstlisting}

\subsection{Factorials, Inverse Factorials \& $nCr$}
Requires prime \code{MOD} and $n<\text{MOD}$ for Fermat inverse.
\begin{lstlisting}
MOD = 10**9 + 7
MAXN = 2000

fact = [1] * (MAXN + 1)
for i in range(1, MAXN + 1):
    fact[i] = fact[i - 1] * i % MOD

inv_fact = [1] * (MAXN + 1)
inv_fact[MAXN] = pow(fact[MAXN], MOD - 2, MOD)
for i in range(MAXN, 0, -1):
    inv_fact[i - 1] = inv_fact[i] * i % MOD

def nCr(n, r):
    if r < 0 or r > n:
        return 0
    return fact[n] * inv_fact[r] % MOD * inv_fact[n-r] % MOD
\end{lstlisting}

% ============================================================
\section{Geometry}

\subsection{Triangle Area --- Heron's Formula}
For sides $a,b,c$ and $s=(a+b+c)/2$:
\[
K=\sqrt{s(s-a)(s-b)(s-c)}.
\]

\subsection{Law of Cosines}
For triangle $ABC$, side $AB$ is opposite angle $C$:
\begin{lstlisting}
from math import acos, degrees
C = degrees(acos((BC*BC + CA*CA - AB*AB)
                 / (2 * BC * CA)))
\end{lstlisting}
Clamp the cosine to $[-1,1]$ if floating-point error is possible.

\subsection{Polygon Area --- Shoelace Formula}
\begin{lstlisting}
def polygon_area2(points):
    # Returns twice the area; exact for integer coordinates.
    area2 = 0
    n = len(points)
    for i in range(n):
        x1, y1 = points[i]
        x2, y2 = points[(i + 1) % n]
        area2 += x1 * y2 - y1 * x2
    return abs(area2)

# area = polygon_area2(points) / 2
\end{lstlisting}

% ============================================================
\section{Graph Algorithms}

\subsection{Iterative DFS with Traversal Depth}
For a tree this is depth from \code{start}; in a general graph it is DFS-tree depth.
\begin{lstlisting}
def dfs_max_depth(start, graph):
    seen = {start}
    stack = [(start, 1)]
    best = 0
    while stack:
        u, dep = stack.pop()
        best = max(best, dep)
        for v in graph[u]:
            if v not in seen:
                seen.add(v)
                stack.append((v, dep + 1))
    return best
\end{lstlisting}

\subsection{Binary Tree Inorder Traversal}
\begin{lstlisting}
def inorder(tree, root):
    ans, st = [], []
    cur = root
    while cur is not None or st:
        while cur is not None:
            st.append(cur)
            cur = tree.get(cur, (None, None))[0]
        cur = st.pop()
        ans.append(cur)
        cur = tree.get(cur, (None, None))[1]
    return ans
\end{lstlisting}

\subsection{N-ary Tree Preorder \& Postorder}
\begin{lstlisting}
def dfs_orders(tree, root):
    pre, post = [], []
    st = [(root, 0)]
    while st:
        u, i = st.pop()
        if i == 0:
            pre.append(u)
        ch = tree.get(u, [])
        if i < len(ch):
            st.append((u, i + 1))
            st.append((ch[i], 0))
        else:
            post.append(u)
    return pre, post
\end{lstlisting}

\subsection{Dijkstra --- Non-negative Edge Weights}
Adjacency list: \code{graph[u] = [(v,w), ...]}.
\begin{lstlisting}
import heapq

def dijkstra(n, graph, src=1):
    INF = float('inf')
    dist = [INF] * (n + 1)
    par = [-1] * (n + 1)
    dist[src] = 0
    pq = [(0, src)]

    while pq:
        du, u = heapq.heappop(pq)
        if du != dist[u]:
            continue
        for v, w in graph[u]:
            nd = du + w
            if nd < dist[v]:
                dist[v] = nd
                par[v] = u
                heapq.heappush(pq, (nd, v))
    return dist, par

def restore_path(par, src, target):
    path = []
    u = target
    while u != -1:
        path.append(u)
        if u == src:
            return path[::-1]
        u = par[u]
    return []
\end{lstlisting}

\subsection{Topological Sort --- Kahn's Algorithm}
\begin{lstlisting}
from collections import deque

def topo_sort(graph, n):
    indeg = [0] * (n + 1)
    for u in range(1, n + 1):
        for v in graph[u]:
            indeg[v] += 1

    q = deque(i for i in range(1, n + 1)
              if indeg[i] == 0)
    order = []
    while q:
        u = q.popleft()
        order.append(u)
        for v in graph[u]:
            indeg[v] -= 1
            if indeg[v] == 0:
                q.append(v)

    return order if len(order) == n else []  # [] => cycle
\end{lstlisting}

% ============================================================
\section{Data Structures}

\subsection{DSU / Union--Find}
1-indexed, iterative path compression, union by size.
\begin{lstlisting}
n = 10
parent = list(range(n + 1))
sz = [1] * (n + 1)

def find(u):
    root = u
    while parent[root] != root:
        root = parent[root]
    while u != root:
        nxt = parent[u]
        parent[u] = root
        u = nxt
    return root

def union(u, v):
    u, v = find(u), find(v)
    if u == v:
        return False
    if sz[u] < sz[v]:
        u, v = v, u
    parent[v] = u
    sz[u] += sz[v]
    return True

# Same component: find(a) == find(b)
# Components after all unions:
# len({find(i) for i in range(1, n + 1)})
\end{lstlisting}

% ============================================================
% FENWICK TREE — V2 UPDATED SECTION
% ============================================================
% Dedicated full-width page with compact code and notes.
% Other sections retain their original code.

\end{multicols}
\clearpage

\begingroup

% Compact code style only for Fenwick Tree
\lstset{
  basicstyle=\ttfamily\fontsize{7.75}{9.3}\selectfont,
  commentstyle=\color{black},
  columns=fullflexible,
  breaklines=true,
  breakatwhitespace=true,
  aboveskip=2pt,
  belowskip=2pt
}

\newcommand{\bithead}[1]{%
  \par\vspace{5pt}\noindent
  {\bfseries\fontsize{9.5}{11}\selectfont #1}%
  \par\vspace{1pt}\hrule\vspace{3pt}
}

\newcommand{\bitcode}[1]{%
  {\ttfamily\fontsize{7.2}{8.5}\selectfont\detokenize{#1}}%
}

\newcommand{\bitpattern}[3]{%
  {\bfseries #1}\par
  {\fontsize{8.4}{9.8}\selectfont
    \textbf{Problem:} #2\par
    \textbf{Approach:} #3\par
  }\vspace{3pt}
}

\subsection{Fenwick Tree / Binary Indexed Tree}

\vspace{2pt}

{\small
\textbf{1-based indexing.}
Build $O(N)$, point update and prefix/range queries
$O(\log N)$, space $O(N)$.
}

\vspace{4pt}

\noindent
\begin{minipage}[t]{0.493\textwidth}
\vspace{0pt}

\bithead{Fenwick Tree / Binary Indexed Tree}

\begin{lstlisting}
N = 200005
bit = [0] * (N + 1)
arr = [0] * (N + 1)

# Build from a 0-indexed list, O(N)
def build(values):
    if len(values) > N:
        raise ValueError("Too many values")
    bit[:] = [0] * (N + 1)
    arr[:] = [0] * (N + 1)
    for i, x in enumerate(values, 1):
        arr[i] = x
        bit[i] = x
    for i in range(1, N + 1):
        j = i + (i & -i)
        if j <= N:
            bit[j] += bit[i]

# Add val at position i
# (does not replace it)
def update(i, val):
    arr[i] += val
    while i <= N:
        bit[i] += val
        i += i & -i

# Prefix sum [1, i]
def query(i):
    ans = 0
    while i > 0:
        ans += bit[i]
        i -= i & -i
    return ans

# Inclusive range sum [l, r]
def range_query(l, r):
    return query(r) - query(l - 1)

# Replace arr[i] with val
def set_value(i, val):
    update(i, val - arr[i])

# First index i with query(i) >= k
# Needs non-negative frequencies
def kth(k):
    if k <= 0 or k > query(N):
        return -1
    i = 0
    step = 1 << (N.bit_length() - 1)
    while step:
        nxt = i + step
        if nxt <= N and bit[nxt] < k:
            k -= bit[nxt]
            i = nxt
        step >>= 1
    return i + 1
\end{lstlisting}

\bithead{Initialize}

\begin{lstlisting}
# Build once from an existing array:
build([2, 4, 1, 5, 3])
# Or start empty and call update(i, x)
\end{lstlisting}

\end{minipage}%
\hfill\vrule width 0.3pt\hfill
\begin{minipage}[t]{0.475\textwidth}
\vspace{0pt}
\raggedright

\bithead{Recognize These Problem Patterns}

\bitpattern{1. Dynamic Range Sum}
{Given an array, support point additions/assignments
and sum queries over $[l,r]$.}
{Store values by position. Use
\bitcode{update(i,x)},
\bitcode{set_value(i,x)}, and
\bitcode{range_query(l,r)}.}

\bitpattern{2. K-th Smallest in a Multiset}
{Insert/delete integer values (including duplicates)
and find the $k$-th smallest.}
{Store frequencies at value indices.
Insert: \bitcode{update(x,1)};
delete: \bitcode{update(x,-1)};
answer: \bitcode{kth(k)}.}

\bitpattern{3. K-th Remaining Element}
{Given an unsorted array, repeatedly remove and
print the $k$-th still-active element
(CSES List Removals).}
{Store 1 for active original positions and 0
for removed ones:
\bitcode{build([1]*n)},
\bitcode{pos=kth(k)},
print \bitcode{a[pos-1]},
then \bitcode{update(pos,-1)}.}

\bitpattern{4. Count Values in a Range}
{Maintain a multiset and count how many
values lie in $[l,r]$.}
{Store frequencies by value and use
\bitcode{range_query(l,r)} to count occurrences.}

\bithead{Quick Usage Examples}

\begin{lstlisting}
# Normal range sum:
build([2, 4, 1, 5, 3])
range_query(2, 4)  # 10
update(3, 5)       # add 5
range_query(2, 4)  # 15
set_value(2, 10)   # replace index 2

# CSES List Removals (fresh tree):
a = [2, 6, 1, 4, 2]
removals = [3, 1, 3, 1, 1]
build([1] * len(a))
ans = []
for k in removals:
    pos = kth(k)
    ans.append(a[pos - 1])
    update(pos, -1)
# ans: [1, 2, 2, 6, 4]
\end{lstlisting}

\bithead{SUM $\rightarrow$ XOR Customization}

{\fontsize{8.4}{9.8}\selectfont
Replace the following operations; the loops stay unchanged.

\par\vspace{2pt}
\renewcommand{\arraystretch}{1.1}

\begin{tabular}{@{}ll@{}}
\hline
\textbf{Function} & \textbf{Change} \\
\hline
\bitcode{build} & \bitcode{bit[j] ^= bit[i]} \\
\bitcode{update} & \bitcode{arr[i] ^= val}, \bitcode{bit[i] ^= val} \\
\bitcode{query} & \bitcode{ans ^= bit[i]} \\
\bitcode{range_query} & \bitcode{query(r) ^ query(l-1)} \\
\bitcode{set_value} & \bitcode{update(i, val ^ arr[i])} \\
\hline
\end{tabular}

\par\vspace{2pt}
\textbf{Note:} \bitcode{kth(k)} is not supported in XOR mode.
}

\bithead{Important Notes}

\begin{itemize}[
  leftmargin=11pt,
  topsep=0pt,
  itemsep=1pt,
  parsep=0pt
]
\item Indices: $1\le i\le N$.
Never call \bitcode{update(0,x)}.

\item \bitcode{update(i,x)} adds;
\bitcode{set_value(i,x)} replaces.

\item \bitcode{kth(k)} finds the first index
whose prefix sum is $\ge k$.
Values must be non-negative.

\item For frequency queries, never delete
an absent value. Compress large/negative
value coordinates.

\item Reset the tree between independent examples.
\end{itemize}

\end{minipage}

\endgroup


% ============================================================
% PAGE 5: BASIC SEGMENT TREE
% ============================================================
\clearpage
\begingroup

\lstset{
  basicstyle=\ttfamily\fontsize{7.75}{9.3}\selectfont,
  commentstyle=\color{black},
  columns=fullflexible,
  breaklines=true,
  breakatwhitespace=true,
  aboveskip=2pt,
  belowskip=2pt
}

\newcommand{\seghead}[1]{%
  \par\vspace{5pt}\noindent
  {\bfseries\fontsize{9.5}{11}\selectfont #1}%
  \par\vspace{1pt}\hrule\vspace{3pt}
}

\newcommand{\segcode}[1]{%
  {\ttfamily\fontsize{7.15}{8.5}\selectfont
   \detokenize{#1}}%
}

\newcommand{\segpattern}[3]{%
  {\bfseries #1}\par
  {\fontsize{8.4}{9.8}\selectfont
    \textbf{Problem:} #2\par
    \textbf{Approach:} #3\par
  }\vspace{3pt}
}


\subsection{Segment Tree --- Basic (Iterative)}
\vspace{2pt}

{\small
\textbf{0-based indexing, inclusive $[l,r]$.}
Build $O(n)$, point update and range query
$O(\log n)$, space $O(n)$.
}

\vspace{4pt}
\noindent

% ============================================================
% LEFT COLUMN: PYTHON TEMPLATE
% ============================================================

\begin{minipage}[t]{\textwidth}
\begin{minipage}[t]{0.493\linewidth}
\vspace{0pt}

\seghead{Segment Tree — Basic (Iterative)}

\begin{lstlisting}
from math import gcd

# Choose before building (default: MIN)
combine = min
IDENTITY = float('inf')

# Build from a 0-indexed list, O(n)
def seg_build(arr):
    n = len(arr)
    size = 1
    while size < n:
        size <<= 1
    tree = [IDENTITY] * (2 * size)
    tree[size:size + n] = arr
    for i in range(size - 1, 0, -1):
        tree[i] = combine(tree[2*i],
                          tree[2*i+1])
    return tree, size

# Inclusive query [l, r]
def seg_query(tree, size, l, r):
    l += size
    r += size
    left = right = IDENTITY
    while l <= r:
        if l & 1:
            left = combine(left, tree[l])
            l += 1
        if not (r & 1):
            right = combine(tree[r], right)
            r -= 1
        l >>= 1
        r >>= 1
    return combine(left, right)

# Replace arr[i] with val
def seg_update(tree, size, i, val):
    i += size
    tree[i] = val
    i >>= 1
    while i:
        tree[i] = combine(tree[2*i],
                          tree[2*i+1])
        i >>= 1

# Add val to arr[i] (point addition)
def seg_add(tree, size, i, val):
    old = tree[size + i]
    seg_update(tree, size, i, old + val)
\end{lstlisting}


\seghead{Initialize}

\begin{lstlisting}
a = [8, 3, 5, 1, 6]
tree, size = seg_build(a)
seg_query(tree, size, 1, 4)  # 1
seg_update(tree, size, 3, 9)
seg_query(tree, size, 1, 4)  # 3
seg_add(tree, size, 1, -2)
seg_query(tree, size, 1, 4)  # 1
\end{lstlisting}

\end{minipage}%
\hfill\vrule width 0.3pt\hfill%
% RIGHT COLUMN: CUSTOMIZATIONS AND APPLICATIONS
\begin{minipage}[t]{0.475\linewidth}
\vspace{0pt}
\raggedright

\seghead{Customize: \segcode{combine} + \segcode{IDENTITY}}

{\fontsize{8.4}{9.8}\selectfont
Change \textbf{both} lines \emph{before}
\segcode{seg_build(arr)}.
The rest of the template stays the same.

\par\vspace{2pt}

\renewcommand{\arraystretch}{1.11}
\begin{tabular}{@{}lll@{}}
\hline
\textbf{Query} & \textbf{combine} & \textbf{IDENTITY} \\
\hline
SUM & \segcode{lambda a,b:a+b} & \segcode{0} \\
MIN & \segcode{min} & \segcode{float('inf')} \\
MAX & \segcode{max} & \segcode{float('-inf')} \\
XOR & \segcode{lambda a,b:a^b} & \segcode{0} \\
GCD & \segcode{gcd} & \segcode{0} \\
AND & \segcode{lambda a,b:a&b} & \segcode{-1} \\
OR & \segcode{lambda a,b:a|b} & \segcode{0} \\
\hline
\end{tabular}
}


\seghead{Recognize These Problem Patterns}

\segpattern{1. Dynamic Range Sum}
{Point assignments/additions; sum of positions $[l,r]$.}
{Use SUM mode, \segcode{seg_update}, \segcode{seg_add},
\segcode{seg_query}. }

\segpattern{2. Range Minimum / Maximum}
{Update individual values; find smallest/largest in $[l,r]$.}
{Use MIN or MAX mode.}


\segpattern{3. Dynamic Range GCD}
{Update one value; calculate GCD over $[l,r]$.}
{Use \segcode{combine = gcd}, \segcode{IDENTITY = 0}.}

\segpattern{4. Range XOR}
{Point updates and XOR over $[l,r]$.}
{Use XOR mode; aggregate neighboring segments with XOR.}

\segpattern{5. Bitwise AND / OR}
{Point updates and bitwise AND/OR over $[l,r]$.}
{Use AND (identity $-1$) or OR (identity $0$).}


\seghead{Quick Usage: SUM Mode}

\begin{lstlisting}
combine = lambda a, b: a + b
IDENTITY = 0
tree, size = seg_build([2, 4, 1, 5, 3])
seg_query(tree, size, 1, 3)  # 10
seg_update(tree, size, 2, 7)
seg_query(tree, size, 1, 3)  # 16
seg_add(tree, size, 1, 3)
seg_query(tree, size, 1, 3)  # 19
\end{lstlisting}


\seghead{Important Notes}

\begin{itemize}[
  leftmargin=11pt,
  topsep=0pt,
  itemsep=1pt,
  parsep=0pt
]

\item Queries are \textbf{inclusive} $[l,r]$;
indices start at $0$. Use a non-empty array
and valid indices.

\item \segcode{seg_update(i,x)}
\textbf{replaces};
\segcode{seg_add(i,x)}
\textbf{adds}, regardless of query mode.
Pass \segcode{tree,size} first.

\item The combine operation must be
\textbf{associative}; identity must be
neutral on both sides.

\item Two query accumulators preserve order
for noncommutative operations.

\item Rebuild the tree after changing
combine/identity.
For range updates use a
\textbf{Lazy Segment Tree}.

\end{itemize}

\end{minipage}%
\end{minipage}

\endgroup


% ============================================================
% PAGE 6: LAZY SEGMENT TREE (ITERATIVE)
% ============================================================
\clearpage
\begingroup
\lstset{
  basicstyle=\ttfamily\fontsize{7.5}{8.45}\selectfont,
  commentstyle=\color{black},
  columns=fullflexible,
  breaklines=false,
  aboveskip=1pt,
  belowskip=1pt
}

\newcommand{\lazyhead}[1]{%
  \par\vspace{5pt}\noindent
  {\bfseries\fontsize{9.2}{11}\selectfont #1}%
  \par\vspace{1pt}\hrule\vspace{3pt}
}
\newcommand{\lazycode}[1]{%
  {\ttfamily\fontsize{7.1}{8.3}\selectfont\detokenize{#1}}%
}
\newcommand{\lazypattern}[3]{%
  {\bfseries #1}\par
  {\fontsize{8.35}{9.7}\selectfont
  \textbf{Problem:} #2\par
  \textbf{Approach:} #3\par}\vspace{3pt}
}

\subsection{Lazy Segment Tree --- Iterative}
\vspace{2pt}
{\small\textbf{0-based, inclusive $[l,r]$.}
Range add/set + SUM/MIN/MAX queries $O(\log n)$;
build $O(n)$, space $O(n)$. No recursion.}
\par\vspace{4pt}
\noindent
\begin{minipage}[t]{\textwidth}
\begin{minipage}[t]{0.493\linewidth}
\vspace{0pt}

\begin{lstlisting}
# SUM / MIN / MAX; 0-based, inclusive [l,r]
MODE = "SUM"

def build(a):
    global n,size,height,tree,lazy_add,lazy_set,seg_len
    global combine,IDENTITY
    n = len(a)
    if not n: raise ValueError("Empty array")
    if MODE == "SUM":
        combine,IDENTITY = (lambda x,y:x+y),0
    elif MODE == "MIN":
        combine,IDENTITY = min,float('inf')
    elif MODE == "MAX":
        combine,IDENTITY = max,float('-inf')
    else: raise ValueError("Invalid MODE")
    size = 1 << (n-1).bit_length()
    height = size.bit_length()-1
    tree = [IDENTITY]*(2*size)
    tree[size:size+n] = a
    lazy_add = [0]*size
    lazy_set = [None]*size
    seg_len = [1]*(2*size)
    for p in range(size-1,0,-1):
        seg_len[p] = seg_len[2*p]+seg_len[2*p+1]
        tree[p] = combine(tree[2*p],tree[2*p+1])

def _set(p,x):
    tree[p] = x*seg_len[p] if MODE=="SUM" else x
    if p < size:
        lazy_set[p],lazy_add[p] = x,0

def _add(p,x):
    tree[p] += x*seg_len[p] if MODE=="SUM" else x
    if p < size:
        if lazy_set[p] is not None: lazy_set[p] += x
        else: lazy_add[p] += x

def _push(p):
    if lazy_set[p] is not None:
        _set(2*p,lazy_set[p]); _set(2*p+1,lazy_set[p])
        lazy_set[p] = None
    if lazy_add[p]:
        _add(2*p,lazy_add[p]); _add(2*p+1,lazy_add[p])
        lazy_add[p] = 0

def _push_paths(L,R):
    for h in range(height,0,-1):
        if L & ((1<<h)-1): _push(L>>h)
        if R & ((1<<h)-1): _push((R-1)>>h)

def _pull_paths(L,R):
    for h in range(1,height+1):
        if L & ((1<<h)-1):
            p = L>>h
            tree[p] = combine(tree[2*p],tree[2*p+1])
        if R & ((1<<h)-1):
            p = (R-1)>>h
            tree[p] = combine(tree[2*p],tree[2*p+1])

def _apply(l,r,x,fn):
    L,R = l+size,r+size+1
    _push_paths(L,R)
    a,b = L,R
    while a < b:
        if a&1: fn(a,x); a += 1
        if b&1: b -= 1; fn(b,x)
        a >>= 1; b >>= 1
    _pull_paths(L,R)

def range_add(l,r,x): _apply(l,r,x,_add)
def range_set(l,r,x): _apply(l,r,x,_set)
def add_value(i,x): range_add(i,i,x)
def set_value(i,x): range_set(i,i,x)

def query(l,r):
    L,R = l+size,r+size+1
    _push_paths(L,R)
    left = right = IDENTITY
    while L < R:
        if L&1: left = combine(left,tree[L]); L += 1
        if R&1: R -= 1; right = combine(tree[R],right)
        L >>= 1; R >>= 1
    return combine(left,right)
\end{lstlisting}

\end{minipage}%
\hfill\vrule width 0.3pt\hfill%
\begin{minipage}[t]{0.475\linewidth}
\vspace{0pt}
\raggedright

\lazyhead{Configuration}

{\fontsize{8.35}{9.7}\selectfont
Set \lazycode{MODE} to \lazycode{"SUM"},
\lazycode{"MIN"}, or \lazycode{"MAX"} before
\lazycode{build(a)}. Rebuilding resets all lazy tags.

\par\vspace{4pt}
\begin{tabular}{@{}lll@{}}
\hline
\textbf{Mode} & \textbf{Combine} & \textbf{Neutral}\\
\hline
SUM & \lazycode{x+y} & 0\\
MIN & \lazycode{min(x,y)} & $+\infty$\\
MAX & \lazycode{max(x,y)} & $-\infty$\\
\hline
\end{tabular}
}

\lazyhead{Update Rules}

{\fontsize{8.35}{9.7}\selectfont
\begin{tabular}{@{}lll@{}}
\hline
\textbf{Mode} & \textbf{Add $x$} & \textbf{Set to $x$}\\
\hline
SUM & \lazycode{tree[p] += x * seg_len[p]} &
      \lazycode{x * seg_len[p]}\\
MIN/MAX & \lazycode{tree[p] += x} &
          \lazycode{x}\\
\hline
\end{tabular}

\par\vspace{3pt}
\textbf{Priority:} SET overrides old additions.
ADD after SET increases the pending set value.
\lazycode{None} means no pending SET; zero is valid.
}

\lazyhead{Recognize These Problem Patterns}

\lazypattern{1. Range Add + SUM}
{Add $x$ to every element in $[l,r]$; query range sums.}
{SUM mode; \lazycode{range_add(l,r,x)}, \lazycode{query(l,r)}.}

\lazypattern{2. Range Set + SUM}
{Assign all values in $[l,r]$ to $x$; query sums.}
{SUM mode; \lazycode{range_set(l,r,x)}.}

\lazypattern{3. Range Add/Set + MIN/MAX}
{Apply updates to intervals; find minimum/maximum.}
{Use MIN or MAX mode; same update/query functions.}

\lazypattern{4. Overlapping Updates}
{Several range additions and assignments overlap.}
{Lazy tags preserve operation order; later SET overwrites
old updates, later ADD increments pending SET.}

\lazyhead{Quick Usage (SUM)}

\begin{lstlisting}
MODE = "SUM"
build([1, 2, 3, 4])
range_add(0, 3, 5)  # [6, 7, 8, 9]
range_set(1, 2, 0)  # [6, 0, 0, 9]
range_add(2, 3, 2)  # [6, 0, 2, 11]
query(0, 3)        # 19
set_value(0, 4)    # replace a[0]
add_value(1, 3)    # add 3 to a[1]
\end{lstlisting}

\lazyhead{Important Notes}

\begin{itemize}[
  leftmargin=11pt,
  topsep=0pt,
  itemsep=2pt,
  parsep=0pt
]

\item Valid range: $0\le l\le r<n$;
endpoints inclusive.

\item Internally use half-open $[L,R)$
with \lazycode{R=r+size+1}.

\item Push boundary paths \emph{before}
applying/querying; pull affected ancestors
\emph{after} updating.

\item Range ADD / SET both $O(\log n)$;
setting a point is just a length-one range update.

\item Choose MODE before build;
do not change it mid-operation.

\end{itemize}

\end{minipage}%
\end{minipage}

\endgroup


\end{document}